\documentclass[11pt]{article}
\usepackage[a4paper,margin=23mm]{geometry}
\usepackage{amsmath,amssymb,array,booktabs}
\usepackage[hidelinks]{hyperref}
\setlength{\parindent}{0pt}
\setlength{\parskip}{6pt}
\newcommand{\C}{\mathbb C}
\newcommand{\End}{\operatorname{End}}
\newcommand{\Span}{\operatorname{span}}
\begin{document}
\begin{center}
{\Large A simple q-Schur cell module with a divisible hook}\\[4pt]
{\large Problem 00000001770: the necessary direction at $q=-1$}\\[3pt]
lizaixi01\qquad 25 October 2026
\end{center}

\textbf{Exact scope.}
The printed conjecture asserts that q-Schur cell modules are semisimple
if and only if the root order $m$ divides none of their hook lengths,
for $m\le5$. We refute only the necessary direction. Over $\C$, take
$m=2$, the genuine primitive root $q=-1$, the algebra $S_q(2,2)$, and
$\lambda=(1,1)$. Its actual right cell module is simple of dimension one,
whereas the top box has hook length two. Both rank and degree equal two;
no $n<r$ boundary convention is used. The separate $m=3,4,5$ table and
hybridization claim is not addressed. Neither is the entire algebra claimed
semisimple.

\textbf{Hecke algebra and the complete permutation sum.}
The rank-two Hecke presentation is
\[
 H=\C[T]/((T+1)^2)\cong\C[\epsilon]/(\epsilon^2),\qquad T=\epsilon-1.
\]
Lean constructs the native \texttt{AdjoinRoot} quotient and inverse algebra
maps to \texttt{TrivSqZeroExt}; the Hecke relation is proved. For every
weak two-composition $\mu$ of two, its actual Young row stabilizer in
$\mathfrak S_2$ defines $m_\mu=\sum_{w\in\mathfrak S_\mu}T_w$ and
$M^\mu=m_\mu H$. There are exactly the three compositions
$(2,0),(1,1),(0,2)$, with sums $\epsilon,1,\epsilon$. Consequently the
complete permutation sum is genuinely isomorphic as an $H$-module to
\[
 M=\C\oplus H\oplus\C,\qquad h\cdot c=h_0c
 \quad(h=h_0+h_1\epsilon).
\]
The file constructs all the actual ideals, their finite direct sum, the
module equivalence, and conjugation onto the native
$\End_H(\bigoplus_\mu M^\mu)$. The coordinate calculation below is
transported to this full endomorphism algebra.

Put $u=(0,1,0)$, $b_0=(1,0,0)$, $b_1=(0,\epsilon,0)$,
$b_2=(0,0,1)$. Let $a_0,a_1,a_2$ extract the first coordinate, the
constant coefficient of the middle coordinate, and the last coordinate.
Define $\phi_{ij}(x)=a_j(x)b_i$ and let $e$ project onto $H$.
Every actual $H$-linear endomorphism $f$ has a unique expansion
\[
 f=\alpha(f)e+\sum_{i,j=0}^2 c_{ij}(f)\phi_{ij},
 \qquad \alpha(f)=\text{constant coefficient of }f(u)_H.
\]
This is a proved native basis, so the true endomorphism dimension is ten.
In particular, $e$ is not the global identity:
$1=e+\phi_{00}+\phi_{22}$.

\newpage
\textbf{Identification of the actual cellular basis.}
All Young diagrams of size two are classified as $(2)$ or $(1,1)$.
Using native semistandard Young tableaux with entries $0,1$, the row shape
has exactly the fillings $00,01,11$; the column shape has only
$\begin{smallmatrix}0\\1\end{smallmatrix}$.
These are the usual entries $1,2$ with a zero-based relabeling. Their
actual contents give the three weak compositions above. For each shape,
the unique standard tableau has content $(1,1)$; its true entry permutation
and every coarsening to the specified content are computed.

The file constructs the Murphy elements and the actual finite sums
$m_{UV}$ over standard tableau pairs with those coarsenings. They evaluate
to $\epsilon$ for the row shape and to $1$ for the column shape. The maps
on \emph{all} permutation generators satisfy the defining formula
\[
 \bigl(\Phi_{UV}(m_\mu h)\bigr)_\nu
 =\begin{cases}
 m_{UV}h,&\mu=\operatorname{type}(V),\ \nu=\operatorname{type}(U),\\
 0,&\text{otherwise}.
 \end{cases}
\]
Thus the nine row maps are precisely $\phi_{ij}$, and the column map is
precisely the regular-summand projection $e$. This is the concrete
rank-two instance of Mathas's formula (4.13), not a labeling assumption.
All tableau pairs form a native basis of the full actual q-Schur algebra.

The multiplication, with $d=(1,0,1)$, is proved directly:
\[
 \phi_{ij}\phi_{kl}=\delta_{jk}d_k\phi_{il},\quad
 e\phi_{ij}=\delta_{i1}\phi_{ij},\quad
 \phi_{ij}e=\delta_{j1}\phi_{ij},\quad e^2=e.
\]
The complex-linear map $\phi_{ij}\mapsto\phi_{ji}$, $e\mapsto e$
is proved involutive and multiplication-reversing. It involves no complex
conjugation. Right multiplication of a row basis element has coefficients
independent of its left tableau. Modulo the row span, right multiplication
of $e$ has coefficient $\alpha(f)$. These are the actual two cellular
triangularity statements.

Dominance is computed from the true row-length prefix sums:
$(2)\triangleright(1,1)$. The actual strictly higher cell span for the
column shape is therefore
\[
 I=\Span_\C\{\phi_{ij}:0\le i,j\le2\}=\ker\alpha.
\]
The equality is derived from the full endomorphism expansion. The map
$\alpha$ is an actual unital algebra homomorphism, so this is the true
two-sided higher ideal. For the row shape the strictly higher span is zero.
The basis, linear anti-involution, dominance spans and their native
endomorphism-algebra transports are all proved in Lean.

\newpage
\textbf{Native cell module and the contradiction.}
The published right cell definition for the column is the cyclic
submodule
\[
 W^{(1,1)}=(e+I)S\ \subseteq S/I.
\]
Lean forms the kernel as a native right submodule, equivalently a left
$S^{\mathrm{op}}$-submodule, and the native module quotient. The first
isomorphism theorem for the actual surjective map $\alpha$ gives
\[
 S/I\ \cong\ \C,\qquad [f]\longmapsto\alpha(f),\qquad
 z\cdot f=z\alpha(f).
\]
Since $\alpha(e)=1$, the cyclic cell submodule is all of this quotient.
Every nonzero vector generates it: for $x\ne0$ and any $y$, the algebra
scalar $(y/x)1$ acts by $y/x$. Hence the native
$S^{\mathrm{op}}$-module is simple and therefore semisimple. Its actual
complex dimension is one, proved by a linear equivalence, not assigned.

These constructions are repeated directly over the native full algebra
$\End_H(\bigoplus_\mu M^\mu)$ using the proved conjugation equivalence:
the transported higher ideal, actual projection, right quotient, cyclic
cell submodule, simplicity and dimension are all present there. Thus the
conclusion is not merely an assertion about an unrelated coordinate
algebra or the underlying complex vector space.

For the genuine column Young diagram, Lean defines the arm and leg as
finite sets of boxes strictly to the right and below. At $(0,0)$ they have
sizes zero and one, so the actual hook length is $0+1+1=2$. Therefore
\[
 W^{(1,1)}\text{ is semisimple},\qquad 2\mid h_{(0,0)}=2.
\]
This negates the printed necessary implication at $m=2\le5$.
The final native-algebra statements are
\texttt{actual\_qSchur\_counterexample} and
\texttt{actual\_necessary\_hook\_condition\_false}
in namespace \texttt{TLMC1770}.

\textbf{Definition boundary and primary source.}
We use Andrew Mathas, \emph{Iwahori-Hecke Algebras and Schur Algebras of
the Symmetric Group}, the actual downloaded notes at
\href{https://web.math.ucsb.edu/~bigelow/books/mathas.pdf}{the UCSB-hosted PDF}:
printed pp.~4,17--20,36,39--42. These pages specify the Hecke relation,
Young sums, permutation modules, semistandard/Murphy maps, formula (4.13),
the linear cellular anti-involution, and the right cell quotient.
The formalization implements this particular complete rank-two model
directly. It does not invoke a general q-Schur classification theorem or
add a Lean axiom identifying an arbitrary algebra or module with this model.

\textbf{Reproduction.}
The package contains the complete bilingual source, this LaTeX/PDF,
the pinned Lean4.33.0/Mathlib project, ordinary exhaustive namespace
axiom checks, original fresh-kernel receipts and actual thread-mask samples.
Run \texttt{lean --version} in \texttt{lean/}, then
\texttt{bash lean/verify.sh} from the package directory for build,
warnings-as-errors axiom checks and the fresh-kernel command:
\begin{center}\texttt{lake env leanchecker --fresh Main}\end{center}
Only \texttt{propext}, \texttt{Classical.choice} and \texttt{Quot.sound}
are allowed. Production runtimes are diagnostic; no controlled efficiency
gain is claimed. Organizer review remains pending.
\end{document}
